\section*{Scope and reading guide}
\addcontentsline{toc}{section}{Scope and reading guide}
\label{p14:comp:guide}

This companion and the associated article, \emph{Sharp Poisson
approximation and second-order correction for long multiplicative runs},
form a single mathematical work. They share the model, notation,
version number, and date. The article follows the exact-marginal
proof of the optimal quadratic approximation. This companion gives
the detailed proofs of the technical inputs and retains the further
results on arithmetic geometry, tails, and deletion.

References beginning with T belong to this companion. Other
numbered references belong to the article. The two PDFs have
cross-document links when kept together under their supplied
filenames. Each cited statement also has a stable printed number,
so the proof can be checked without those links.

\paragraph{Common model and conventions.}
The completely multiplicative function \(f\) has independent fair
signs at primes. Both signs of a constant run are counted. With
\[
 I_N=[N,2N)\cap\Z,\quad |L-\log_2N|\le C,\quad
 B=L+1,\quad p=2^{-L},\quad \lambda_N=Np,
\]
the event at \(x\in I_N\) and its count are
\[
 J_x=\ind_{\{f(x-1)=-f(x),\ f(x)=\cdots=f(x+L-1)\}},
 \qquad Z_N=\sum_{x\in I_N}J_x.
\]
We write \(p_x=\E J_x\), \(\mu_N=\E Z_N\),
\(\alpha_N=\Var Z_N-\mu_N\), \(\Pi_\mu=\Pois(\mu)\),
and \(\Delta=\delta_1-\delta_0\). Products of measures
are convolutions, \(\|\nu\|_1=\sum_k|\nu\{k\}|\),
and total variation for probabilities is half this norm.
All unmarked logarithms are natural. Constants may depend on
the fixed critical-band width \(C\) and on parameters explicitly
declared fixed.

The space \(S_x\) is the image, under prime valuations modulo two,
of the even selections in \(V_x=\{x-1,\ldots,x+L-1\}\).
The exact formulation uses all \(x\) with \(p_x>0\), dimensions
\(d_x=\dim S_x\), and the relative nullity
\(\sum_{x\in I}d_x-\dim\sum_{x\in I}S_x\).
The regular-set argument restricts to \(G_N\), where \(d_x=L\)
and \(p_x=p\). Its homogeneous notation and its restricted count
are always identified when that argument is used. The raw arithmetic
nullity \(\rho(I)=|I|L-\dim\sum_{x\in I}S_x\) includes local
kernels; it agrees with the effective nullity on \(G_N\).
Section~\ref{p13:transfer:section} explicitly uses \(\rho\)
as a local abbreviation for the effective nullity.

\paragraph{Technical inputs for the optimal scalar theorem.}
The arithmetic estimates in Sections~\ref{p1:arith:appendix}
and~\ref{p1:arith:supplement} supply the split-product bounds,
variable-coefficient counts, and channel budgets. The proofs
associated with Sections~\ref{p1:fixed:section},
\ref{p13:short:section}, and~\ref{p13:half:section} establish
fixed-source propagation and conditional half counts. The
rigidity and triple-counting proofs then supply the inputs to
Sections~\ref{p13:rig:section} and~\ref{p13:third:section}.
The article states the required hypotheses at each interface
and retains the reconstruction, its separately uniform tail,
and the quartic lower-bound argument.

\paragraph{Entering a proof and returning to the article.}
The following local routes lead to the complete proof subsection,
then to the next use in the article. A subsection and a lemma can
have the same printed number; the middle column names the
subsection and links to its heading.

\begin{center}
\small
\begin{tabularx}{\textwidth}{@{}>{\raggedright\arraybackslash}p{0.23\textwidth}>{\raggedright\arraybackslash}p{0.40\textwidth}>{\raggedright\arraybackslash}X@{}}
\toprule
Entry in the article & Proof in this companion & Return to the article\\
\midrule
Proposition~\ref{p13:short:moments}
& Subsection~\ref{p14:comp:short:moments}: uniform moments, using
the prescribed-divisor estimates above it
& Corollary~\ref{p13:short:old-space}: short projections into the old space\\[4pt]
Theorem~\ref{p13:half:conditional}
& Subsection~\ref{p14:comp:half:iteration}: threshold iteration,
using the certificate count of Lemma~\ref{p13:half:saturation}
& Corollary~\ref{p13:half:connected}: connected hosts at fixed order\\[4pt]
Theorem~\ref{p13:rig:common-curve}
& Subsection~\ref{p14:comp:rigidity:common-curve}: one curve for
the whole relation space
& Proposition~\ref{p13:rig:grid-fibres}: integer fibres, followed
by the degree-uniform point count\\[4pt]
Proposition~\ref{p13:third:aligned-family}
& Subsection~\ref{p14:comp:third:aligned-family}: all triples above
the fixed channels, including exceptional pairs
& Corollary~\ref{p13:third:aligned}: strong aligned pairs\\[4pt]
Lemma~\ref{p13:transfer:dense-targets}
& Subsection~\ref{p14:comp:transfer-dense-targets}: two dense targets
over prescribed sources
& Lemma~\ref{p13:transfer:compressed-target}, then the primitive-mass estimates\\
\bottomrule
\end{tabularx}
\end{center}

\paragraph{The regular-set reconstruction and positional results.}
The complete earlier reconstruction is retained in
Sections~\ref{p1:core:section}--\ref{p1:close:section} and
the associated stopping proofs. Its weighted primitive and
compressed-pair estimates support the unsigned positional
comparison and the deletion arguments in
Section~\ref{p1:limits:positions}. This branch is needed for
those results; the scalar optimal proof is assembled directly
from its stated exact-marginal inputs. Any proof below that
invokes Theorem~\ref{p1:thm:main} is a consequence of that
theorem, not an input to it.

\paragraph{Further consequences and limitations.}
Sections~\ref{p1:offset:section}, \ref{p1:consequences:section},
and~\ref{p1:uniform:appendix} give refinements of the offset
geometry, interval dependence, and uniform counting estimates.
Section~\ref{p1:limits:section} treats exact local blocks,
effective tails, modified counts, Palm identities, and the cost
of positive forcing. Section~\ref{p1:moment-obstruction}
explains why fixed moments alone cannot imply a polynomial
rate for a complete law. These statements retain their own
quantifiers; none enlarges a fixed-order estimate to a growing
range without a separate proof.

\paragraph{Which final sections enter an article result?}
The dependency map in Section~\ref{p1:intro:dependencies} identifies
the independent fixed-moment proof, the triple bound, the
reconstruction, and the quartic lower bound. Within the final five
sections of this companion, the roles are as follows.

\begin{center}
\small
\begin{tabularx}{\textwidth}{@{}>{\raggedright\arraybackslash}p{0.32\textwidth}>{\raggedright\arraybackslash}X@{}}
\toprule
Part of the companion & Role in the associated work\\
\midrule
Positional comparison in
Section~\ref{p1:limits:positions}
& Supplies the proof of Lemma~\ref{p1:limits:vector-reference}
in the article. It uses the regular-set weighted reconstruction;
the scalar optimal theorem does not use it.\\[4pt]
Other parts of Section~\ref{p1:limits:section}
& Local-block, tail, deletion, Palm, and forcing results.
They are additional conclusions with their own assumptions;
they are not inputs to the scalar optimal theorem.\\[4pt]
Sections~\ref{p1:offset:section},
\ref{p1:consequences:section}, \ref{p1:uniform:appendix},
and~\ref{p1:moment-obstruction}
& Further geometry, interval dependence, counting refinements,
and a moment obstruction. The article cites selected statements
to explain scope; its scalar reconstruction does not depend on them.\\
\bottomrule
\end{tabularx}
\end{center}

Here ``additional'' refers to the logical role of a result, not
to an assertion that it forms a separate paper. The quantized
mean identity in Proposition~\ref{p13:optimal:nominal} is proved
directly from the exact marginals; it does not require the
supplementary local-block asymptotics.

The preserved version~1.6 source archive provides
\texttt{python3 check\_all.py} for its historical builds and finite
diagnostics. Version~1.7 is rebuilt from the associated repository with
\texttt{python3 scripts/build\_manuscript\_v17.py}.
Source integrity and statement correspondence are checked by
\texttt{scripts/check\_revised\_sources.py}; finite diagnostics
do not test the uniform asymptotic estimates.
